% ===== Block B14: 原文 p.131–140 | 第 6 章 顶点算子与 KZ 方程（§6.1–6.2） | 译手 tr-ch1c =====
\subsection{黎曼球面 $\PP^1$ 上的共形块}
% [原文 p.131]
本节我们研究当底曲线为黎曼球面 $\PP^1$ 时的真空空间。把 $\PP^1$ 视为 $\C\cup\{\infty\}$，并取 $z$ 为 $\C$ 的整体坐标。对正整数 $N\ge 1$，在 $\PP^1$ 上任取 $N$ 个点 $z_1,z_2,\dots,z_N$，并置
\[
\xi_j=\begin{cases}
z-z_j & \text{若 } z_j\ne\infty,\\
1/z & \text{若 } z_j=\infty.
\end{cases}
\]
于是 $\mathfrak{X}=(\PP^1;z_1,\dots,z_N;\xi_1,\dots,\xi_N)$ 是一条带形式坐标的 $N$-带点曲线。取 $\vec\lambda=(\lambda_1,\dots,\lambda_N)\in(P_\ell)^N$，并置
\[
V_{\vec\lambda}=V_{\lambda_1}\otimes_{\C}\cdots\otimes_{\C}V_{\lambda_N}.
\]
% [Tutor Note: 这里的 $\mathfrak{X}$ 就是 §1.4 意义下的 $N$-带点曲线，只是把底曲线取成 $\PP^1$、把形式坐标 $\xi_j$ 显式写出。核对办法：当 $z_j\ne\infty$ 时 $\xi_j=z-z_j$ 正是 $z_j$ 处的局部坐标（在 $z_j$ 处为零点、在 $\infty$ 处有极点）；当 $z_j=\infty$ 时 $\xi_j=1/z$ 是 $\infty$ 处的局部坐标。后文所有 $\mathrm{Res}_{\xi_j=0}$ 都相对这组坐标取。]
% [Tutor Note: $P_\ell$ 是水平（level）为 $\ell$ 的可积最高权之集，即第 2 章中的整支配权集；$(\theta,\lambda_k)\in\Z_{\ge 0}$ 是 $\lambda_k$ 与最高根 $\theta$ 的配对，下文的 $\ell_k$ 依赖它取整。]
\begin{proposition}[6.1]
自然的限制映射
\[
\Hom_{\C}(\mathcal{H}_{\vec\lambda},\C)\to\Hom_{\C}(V_{\vec\lambda},\C)
\]
诱导出单射同态
\[
j:\mathcal{V}^{\dagger}_{\vec\lambda}(\mathfrak{X})\hookrightarrow\Hom_{\mathfrak{g}}(V_{\vec\lambda},\C).
\]
此外，若没有 $z_i$ 等于 $\infty$，则元素 $\langle\psi|\in\Hom_{\mathfrak{g}}(V_{\vec\lambda},\C)$ 落在 $j$ 的像中，当且仅当对每个 $k=1,2,\dots,N$ 以及任意元素
\[
|\phi_k\rangle=|v_1\otimes\cdots\otimes v_{k-1}\otimes\lambda_k\otimes v_{k+1}\otimes\cdots\otimes v_N\rangle\in V_{\vec\lambda},
\]
都有等式
\begin{equation}
\sum_{\vec m_k,\,|\vec m_k|=\ell_k}\prod_{j\ne k}\binom{\ell_k}{m_j}(z_j-z_k)^{-m_j}\langle\psi|\prod_{i\ne k}\rho_i(X_\theta)^{m_i}|\phi_k\rangle=0,\tag{6.1}
\end{equation}
其中
\[
\ell_k=\ell-(\theta,\lambda_k)+1,
\]
\[
\vec m_k=(m_1,\dots,m_{k-1},m_{k+1},\dots,m_N),\qquad \Z\ni m_i\ge 0,
\]
\[
|\vec m_k|=\sum_{i\ne k}m_i,
\]
且 $|\lambda_k\rangle$ 是 $V_{\lambda_k}$ 的最高权向量。
\end{proposition}
% [Tutor Note: (6.1) 里两处下标不能混：$(z_j-z_k)^{-m_j}$ 与二项式系数 $\binom{\ell_k}{m_j}$ 都取 $j$（即 $j\ne k$ 的那一路），而 $\rho_i(X_\theta)^{m_i}$ 取 $i$；两个乘积记号 $\prod_{j\ne k}$、$\prod_{i\ne k}$ 都跑遍除 $k$ 以外的全部指标。取 $\vec m_k=0$ 项可见 (6.1) 的左端含 $\langle\psi|\phi_k\rangle$。]
% [Tutor Note: $\ell_k=\ell-(\theta,\lambda_k)+1$ 的来历：在 $V_{\lambda_k}$ 中 $X_\theta(-1)$ 作用 $\ell_k$ 次后给出零（见 (2.8) 型的 $\mathfrak{sl}_2$ 关系），故只有 $|\vec m_k|\le\ell_k$ 的项可能非零；核对时用 $(\theta,\theta)=2$ 及 $(\theta,\lambda_k)$ 为整数即可。]
注意，等式 (6.1) 等价于
\[
\langle\Psi|\rho_k(X_\theta(-1))^{\ell_k}|\phi\rangle=0,
\]
只要置 $\langle\psi|=j(\langle\Psi|)$、$\langle\Psi|\in\mathcal{V}^{\dagger}_{\vec\lambda}(\mathfrak{X})$（见下面的 (6.2)）。
% [原文 p.132]
\noindent\textit{证明.}\quad 由规范条件，对 $\langle\Psi|\in\mathcal{V}^{\dagger}_{\vec\lambda}(\mathfrak{X})$ 与 $X\in\mathfrak{g}$ 有
\[
\sum_{j=1}^{N}\langle\Psi|\rho_j(X\otimes 1)=0.
\]
% [Tutor Note: 这条“规范条件”就是总荷为零：$\sum_j\rho_j(X\otimes 1)$ 是 $X$ 的总作用，它必须把真空态打成零。第 5 章用 $\hat{\mathfrak{g}}$ 作用在 $\mathcal{H}_{\vec\lambda}$ 上的不变性来刻画真空空间，这里把它逐分量写出。]
这说明 $j$ 的像落在 $\Hom_{\mathfrak{g}}(V_{\vec\lambda},\C)$ 中。
现在证明 (6.1) 是 $\langle\psi|$ 属于 $j(\mathcal{V}^{\dagger}_{\vec\lambda}(\mathfrak{X}))$ 的必要条件。对 $\langle\psi|=j(\langle\Psi|)$ 与 $h=1/(z-z_k)$，用与定理 3.15 的证明类似的论证，可得如下等式。
\begin{align}
\langle\Psi|\rho_k(X_\theta(-1))^{\ell_k}|\phi\rangle
&=\langle\Psi|\rho_k(X_\theta\otimes h)^{\ell_k}|\phi\rangle\nonumber\\
&=(-1)^{\ell_k}\sum_{\vec m_k,\,|\vec m_k|=\ell_k}\binom{\ell_k}{m_j}(z_j-z_k)^{-m_k}\langle\Psi|\prod_{j\ne k}\rho_j(X_\theta)^{m_j}|\phi\rangle.\tag{6.2}
\end{align}
% [Tutor Note: (6.2) 的推导要点：$X_\theta(-1)$ 是 $X_\theta\otimes z^{-1}$ 的 $-1$ 次分量，而 $h=1/(z-z_k)$ 在 $z=z_k$ 附近展开为 $\sum_{m\ge0}(z-z_k)^{m-1}z_k^{-m-1}$ 一类的级数；把它代入 $\rho_k(X_\theta\otimes h)^{\ell_k}$ 后，凡 $h$ 在 $z_j$（$j\ne k$）处全纯的项都可移到别的张量因子上，剩下的极点项给出 $(z_j-z_k)^{-1}$ 型因子。$(-1)^{\ell_k}$ 来自 $\rho$ 的斜对称性（把 $\ell_k$ 个 $X_\theta$ 逐个移过）。核对办法：取 $\ell_k=1$、$\vec m_k=0$ 退化为 $\langle\Psi|\rho_k(X_\theta\otimes h)|\phi\rangle$。]
此外，由定理 2.4 有
\[
X_\theta(-1)^{\ell_k}|\lambda_k\rangle=0.
\]
因此等式 (6.2) 蕴含等式 (6.1)。
下面证明 $j$ 是单射。设 $j(\langle\Psi|)=0$。对 $p$ 作归纳，证明 $\langle\Psi|$ 在 $F_p\mathcal{H}_{\vec\lambda}$ 上的限制 $\langle\Psi|\big|_{F_p(\mathcal{H}_{\vec\lambda})}$ 为零。$p=0$ 时由假设 $j(\langle\Psi|)=0$ 即知。设 $\langle\Psi|\big|_{F_p(\mathcal{H}_{\vec\lambda})}=0$。元素 $|\phi\rangle\in F_{p+1}\mathcal{H}_{\vec\lambda}$ 可写成
\[
\rho_j(X(-n))|v\rangle,\qquad |v\rangle\in F_p\mathcal{H}_{\vec\lambda},
\]
对某个 $j$、正整数 $n$ 与 $X\in\mathfrak{g}$ 成立。置 $f=1/(z-z_j)^n$。于是有
\begin{align*}
\langle\Psi|\phi\rangle&=\langle\Psi|\rho_j(X(-n)|v\rangle\\
&=\langle\Psi|\rho_j(X\otimes f)|v\rangle\\
&=-\sum_{i\ne j}\langle\Psi|\rho_i(X\otimes f)|v\rangle=0,
\end{align*}
因为 $f$ 在 $z_i$（$i\ne j$）处全纯，故 $\rho_i(X\otimes f)|v\rangle\in F_q\mathcal{H}_{\vec\lambda}$。这就证明了 $\langle\Psi|=0$。
% [Tutor Note: 原文此处把 $F_p$ 印成 $F_q$（下一段归纳里又改用 $q$ 作归纳下标），按上下文应为 $F_p$：$|v\rangle\in F_p$ 而 $\rho_i(X\otimes f)$ 不提高滤过度。又第一行 $\rho_j(X(-n)|v\rangle$ 少一个右括号，正确写法是 $\rho_j(X(-n))|v\rangle$。两处均照排，见报告 ERRATA。]
最后证明条件 (6.1) 的充分性。先证明 $\Hom_{\mathfrak{g}}(V_{\vec\lambda},\C)$ 中每个元素 $\langle\psi|$ 都可以延拓成元素
\[
\langle\Psi^{(p)}|\in\Hom_{\C}(F_p\mathcal{M}_{\vec\lambda},\C),
\]
其中
\begin{align*}
\mathcal{M}_{\vec\lambda}&=\mathcal{M}_{\lambda_1}\otimes_{\C}\cdots\otimes_{\C}\mathcal{M}_{\lambda_N},\\
F_p\mathcal{M}_{\vec\lambda}&=\bigoplus_{p_1+\cdots+p_N=p}F_{p_1}\mathcal{M}_{\lambda_1}\otimes_{\C}\cdots\otimes_{\C}F_{p_N}\mathcal{M}_{\lambda_N},
\end{align*}
而 $\mathcal{M}_{\lambda_j}$ 是以 $\lambda_j$ 为最高权的 Verma 模（(2.7)）。更确切地说，对 $p$ 作归纳，我们证明存在延拓
\[
\langle\Psi^{(p)}|\in\Hom(F_p\mathcal{M}_{\vec\lambda},\C),
\]
使得若
对元素 $f\in H^0(\PP^1,\mathcal{O}_{\PP^1}(*\sum_{j=1}^{N}z_j))$ 与 $|v\rangle\in F_p\mathcal{M}_{\vec\lambda}$ 有
\[
\rho_j(X\otimes f)|v\rangle\in F_p\mathcal{M}_{\vec\lambda},\qquad j=1,2,\dots,N,
\]
则有
\begin{equation}
\sum_{j=1}^{N}\langle\Psi^{(p)}|\rho_j(X\otimes f)|v\rangle=0.\tag{6.3}
\end{equation}
% [Tutor Note: $\mathcal{M}_{\vec\lambda}$ 是 Verma 模的张量积，$F_p$ 按“各因子滤过度之和 $\le p$”定义；(6.3) 的归纳构造之所以可行，是因为 $\Hom_{\C}(F_p\mathcal{M}_{\vec\lambda},\C)$ 有限维、且 $F_{q+1}/F_q$ 由 $\rho_j(-n)|v\rangle$ 型的向量生成（见下面两段）。]
% [原文 p.133]
$p=0$ 时置
\[
\langle\Psi^{(0)}|=\langle\psi|.
\]
于是对常数函数 $f$，等式 (6.3) 成立。今设已求得 $\langle\Psi^{(q)}|$。$F_{q+1}\mathcal{M}(\vec\lambda)$ 的元素 $|\phi\rangle$ 可表为
\[
\rho_j(-n)|v\rangle,
\]
其中 $j$ 适当、$n$ 是正整数、$|v\rangle\in F_q\mathcal{M}_{\vec\lambda}$。置
\[
h(z)=\frac{1}{(z-z_j)^n},
\]
并如下定义 $\langle\Psi^{(q+1)}|$：
\begin{equation}
\langle\Psi^{(q+1)}|\phi\rangle=-\sum_{i\ne j}\langle\Psi^{(q)}|\rho_i(X\otimes h)|v\rangle.\tag{6.4}
\end{equation}
由于 $h(z)$ 在 $z_i$（$i\ne j$）处全纯，故 $\rho_i(X\otimes h)|v\rangle\in F_q\mathcal{M}_{\vec\lambda}$。因此 (6.4) 的右端是良定义的。等式 (6.3) 是定义 (6.4) 与归纳假设的推论。于是得到
\[
\langle\Psi|\in\Hom_{\C}(\mathcal{M}_{\vec\lambda},\C),
\]
它满足规范条件
\begin{equation}
\sum_{j=1}^{N}\langle\Psi|\rho_j(X\otimes f)=0\tag{6.5}
\end{equation}
对一切 $f\in H^0(\PP^1,\mathcal{O}_{\PP^1}(*\sum_{j=1}^{N}z_j))$ 成立。
现在证明 $\langle\Psi|$ 实际上定义了 $\Hom_{\C}(\mathcal{H}_{\vec\lambda},\C)$ 中的一个元素，从而由规范条件 (6.5) 它是 $\mathcal{V}^{\dagger}_{\vec\lambda}(\mathfrak{X})$ 的元素。只需证明对每个 $k$，$1\le k\le N$，有
\begin{equation}
\langle\Psi|(\mathcal{M}_{\lambda_1}\otimes\cdots\otimes\mathcal{M}_{\lambda_{k-1}}\otimes\mathcal{J}_{\lambda_k}\otimes\mathcal{M}_{\lambda_{k+1}}\otimes\cdots\otimes\mathcal{M}_{\lambda_N})=0.\tag{6.6}
\end{equation}
注意 $\mathcal{J}_{\lambda_k}$ 由
\[
U(\hat{\mathfrak{p}}_-)|J_{\lambda_k}\rangle,\qquad |J_{\lambda_k}\rangle=X_\theta(-1)^{\ell-(\theta,\lambda_k)+1}|\lambda_k\rangle
\]
定义，其中 $U(\hat{\mathfrak{p}}_-)$ 是李代数 $\hat{\mathfrak{p}}_-$ 的泛包络代数。于是 $\mathcal{J}_{\lambda_k}$ 是 $\mathcal{M}_{\lambda_k}$ 的极大真 $\hat{\mathfrak{g}}$-子模，且 $\mathcal{H}_{\lambda_k}=\mathcal{M}_{\lambda_k}/\mathcal{J}_{\lambda_k}$。先证明
\begin{equation}
\langle\Psi|(V_{\lambda_1}\otimes\cdots\otimes V_{\lambda_{k-1}}\otimes\mathcal{J}_{\lambda_k}\otimes V_{\lambda_{k+1}}\otimes\cdots\otimes V_{\lambda_N})=0\tag{6.7}
\end{equation}
对滤过次数 $F_\bullet\mathcal{J}_{\lambda_k}:=F_\bullet\mathcal{M}_{\lambda_k}\cap\mathcal{J}_{\lambda_k}$ 作归纳。由上一段的假设有
\[
\langle\Psi|(V_{\lambda_1}\otimes\cdots\otimes V_{\lambda_{k-1}}\otimes J_{\lambda_k}\otimes V_{\lambda_{k+1}}\otimes\cdots\otimes V_{\lambda_N})=0.
\]
设 $\langle\Psi|$ 在 $F_p\mathcal{J}_{\lambda_k}$ 上为零。则 $F_{p+1}\mathcal{J}_{\lambda_k}$ 的元素是如下元素的 $\C$-线性组合：
\[
|v_1\otimes\cdots\otimes v_{k-1}\otimes Y(-n)w\otimes v_{k+1}\otimes\cdots\otimes v_N\rangle,
\]
% [原文 p.134]
其中 $v_i\in V_{\lambda_i}$、$Y\in\mathfrak{g}$、$w\in F_p\mathcal{J}_{\lambda_k}$，$n$ 是正整数。由于 $1/(\zeta-z_k)^n$ 在 $z_i$（$i\ne k$）处全纯，故有
\begin{equation}
Y\otimes\frac{1}{(\zeta-z_k)^n}|v_i\rangle\in V_{\lambda_i}.\tag{6.8}
\end{equation}
于是有
\begin{align*}
&\langle\Psi|v_1\otimes\cdots\otimes v_{k-1}\otimes Y(-n)w\otimes v_{k+1}\otimes\cdots\otimes v_N\rangle\\
&\quad=\Big\langle\Psi\Big|\rho_k\Big(Y\otimes\frac{1}{(\zeta-z_k)^n}\Big)v_1\otimes\cdots\otimes v_{k-1}\otimes w\otimes v_{k+1}\otimes\cdots\otimes v_N\Big\rangle\\
&\quad=-\sum_{i\ne k}\Big\langle\Psi\Big|\rho_i\Big(Y\otimes\frac{1}{(\zeta-z_k)^n}\Big)v_1\otimes\cdots\otimes v_{k-1}\otimes w\otimes v_{k+1}\otimes\cdots\otimes v_N\Big\rangle\\
&\qquad\qquad\qquad\qquad\qquad\qquad\qquad\qquad\qquad\qquad\qquad\qquad\text{（由 (6.5)）}\\
&\quad=0\qquad\text{（由 (6.8) 与归纳假设）}.
\end{align*}
% [Tutor Note: (6.8) 的意思是 $Y\otimes(\zeta-z_k)^{-n}$ 可合法地作用在 $V_{\lambda_i}$ 上：$\zeta$ 是第 $i$ 个因子的形式坐标，而 $(\zeta-z_k)^{-n}$ 在 $z_i$（$i\ne k$）处全纯，故属于第 4 章引理（(4.18) 型）所允许的系数环。核对办法：把 $|v_i\rangle$ 展开成 $\zeta$ 的幂级数，逐项作用即可。]
这样便得到所要的 (6.6)。
其次考虑
\[
\langle\Psi|Y(-n)v_1\otimes\cdots\otimes v_{k-1}\otimes\mathcal{J}_{\lambda_k}\otimes v_{k+1}\otimes\cdots\otimes v_N\rangle,
\]
其中 $Y\in\mathfrak{g}$、$|v_i\rangle\in V_{\lambda_i}$，$n$ 是正整数。由于
\[
Y\otimes\frac{1}{(\zeta-z_k)^n}|v_i\rangle\in V_{\lambda_i},\qquad Y\otimes\frac{1}{(\zeta-z_k)^n}\mathcal{J}_{\lambda_k}\subset\mathcal{J}_{\lambda_k}
\]
对 $i\ge 2$ 成立，由 (6.5) 有
\begin{align}
&\langle\Psi|Y(-n)v_1\otimes\cdots\otimes v_{k-1}\otimes\mathcal{J}_{\lambda_k}\otimes v_{k+1}\otimes\cdots\otimes v_N\rangle\nonumber\\
&\quad=-\sum_{i=2}^{N}\Big\langle\Psi\Big|\rho_i\Big(Y\otimes\frac{1}{(\zeta-z_k)^n}\Big)v_1\otimes\cdots\otimes v_{k-1}\otimes\mathcal{J}_{\lambda_k}\otimes v_{k+1}\otimes\cdots\otimes v_N\Big\rangle\tag{6.9}\\
&\quad=0.\nonumber
\end{align}
如此，对滤过次数
\[
F_\bullet(\mathcal{M}_{\lambda_1}\otimes\cdots\otimes\mathcal{M}_{\lambda_{k-1}}\otimes\mathcal{M}_{\lambda_{k+1}}\otimes\cdots\otimes\mathcal{M}_{\lambda_N})
\]
作归纳，便可证明 (6.6) 成立。\hfill$\square$
\begin{corollary}[6.2]
(1) 当 $N=1$ 时，有
\[
\mathcal{V}^{\dagger}_{\lambda}(\mathfrak{X})=0\qquad(\lambda\ne 0).
\]
(2) 当 $N=2$ 时，有
\[
\mathcal{V}^{\dagger}_{\lambda\mu}(\mathfrak{X})=\begin{cases}\C & \mu=\lambda^{\dagger},\\ 0 & \text{其他情形}.\end{cases}
\]
(3) 当 $N=3$ 且 $\vec\lambda=(\mu,\nu,\lambda)$ 时，有
\[
\mathcal{V}^{\dagger}_{\vec\lambda}(\mathfrak{X})\simeq W_{\vec\lambda},
\]
其中
\[
W_{\vec\lambda}=\{\psi\in\Hom_{\mathfrak{g}}(V_{\vec\lambda},\C)\mid \text{条件 }(*)\},
\]
而条件 $(*)$ 如下给出。设
\[
\mathfrak{k}_\theta=\C X_\theta\oplus\C X_{-\theta}\oplus\C[X_\theta,X_{-\theta}]
\]
是 $\mathfrak{g}$ 的主三维子代数，并设
\[
V_{\lambda}=\bigoplus_{i=0}^{\ell/2}W_{\lambda,i}
\]
是作为 $\mathfrak{k}_\theta$-模的不可约分解。则条件 $(*)$ 是
\[
(*)\qquad \psi|_{W_{\mu,h}\otimes W_{\nu,i}\otimes W_{\lambda,j}}=0,\qquad \text{若 } h+i+j>\ell.
\]
\end{corollary}
% [Tutor Note: 推论 6.2 (2) 里的 $\lambda^{\dagger}$ 是第 5 章的对偶权（$\lambda^{\dagger}=-w_0\lambda$）；(3) 中的 $W_{\lambda,i}$ 是 $\mathfrak{k}_\theta\simeq\mathfrak{sl}_2$ 的（$2i+1$ 维）不可约成分，$h+i+j>\ell$ 时三重张量积的最高权超过水平 $\ell$，故只能取零——这是 $\mathfrak{sl}_2$ 的 Clebsch--Gordan 判定，可用来核对条件 $(*)$ 的阈值。]
% [原文 p.135]
下面考虑关联函数。由于 $\PP^1$ 上没有全纯 1-形式，由定理 3.15 第 5) 条，对 $X\in\mathfrak{g}$ 与 $|v\rangle\in V_{\vec\lambda}$ 有
\begin{equation}
\langle\Psi|X(z)|v\rangle dz=\sum_{j=1}^{N}\frac{1}{z-z_j}\langle\Psi|\rho_j(X)v\rangle dz.\tag{6.10}
\end{equation}
类似地，对 $X,Y\in\mathfrak{g}$ 有
\begin{align}
\langle\Psi|X(z)Y(w)|v\rangle dzdw&=\frac{\ell\cdot(X,Y)}{(z-w)^2}\langle\Psi|v\rangle dzdw\nonumber\\
&\quad+\frac{1}{z-w}\langle\Psi|[X,Y](w)|v\rangle dzdw\nonumber\\
&\quad+\sum_{j=1}^{N}\frac{1}{z-z_j}\langle\Psi|Y(w)|\rho_j(X)v\rangle dzdw\nonumber\\
&\quad+\sum_{j=1}^{N}\frac{1}{z-z_j}\langle\Psi|X(z)|\rho_j(Y)v\rangle dzdw\tag{6.11}
\end{align}
% [Tutor Note: (6.11) 的四项分别是：$T$ 的二次极点项（中心荷/水平 $\ell$ 乘 Killing 型）、$T$ 的一次极点项给的 $[X,Y]$ 交换子、以及两个“与插入点 $z_j$ 的碰撞”项。(6.10) 说 $\langle\Psi|X(z)|v\rangle dz$ 只能是简单极点之和，物理上就是“带点球面上无全纯 1-形式”的陈述；核对办法：取 $N=0$ 时 (6.10) 右端为零，即 $\langle\Psi|X(z)|v\rangle dz$ 必为零。]
由于 $\PP^1\times\PP^1$ 上没有全纯 2-形式，由 (6.10) 可把等式 (6.11) 改写成
\begin{align}
\langle\Psi|X(z)Y(w)|v\rangle dzdw&=\frac{\ell\cdot(X,Y)}{(z-w)^2}\langle\Psi|v\rangle dzdw+\frac{1}{z-w}\sum_{j=1}^{N}\frac{1}{z-z_j}\langle\Psi|\rho_j([X,Y])v\rangle dzdw\nonumber\\
&\quad+2\sum_{j=1}^{N}\sum_{k=1}^{N}\frac{1}{(z-z_j)(w-z_k)}\langle\Psi|\rho_j(X)\rho_k(Y)v\rangle dzdw.\tag{6.12}
\end{align}
这也可写成如下形式。
\begin{align}
\langle\Psi|X(z)Y(w)|v\rangle dzdw&=\frac{\ell\cdot(X,Y)}{(z-w)^2}\langle\Psi|v\rangle dzdw+\frac{1}{w-z}\sum_{j=1}^{N}\frac{1}{w-z_j}\langle\Psi|\rho_j([Y,X])v\rangle dzdw\nonumber\\
&\quad+2\sum_{j=1}^{N}\sum_{k=1}^{N}\frac{1}{(z-z_k)(w-z_j)}\langle\Psi|\rho_j(Y)\rho_k(X)v\rangle dzdw.\tag{6.13}
\end{align}
% [Tutor Note: (6.12)$\to$(6.13) 只是把 $z\leftrightarrow w$、$X\leftrightarrow Y$ 后重命名：交换后 $[X,Y]\to[Y,X]$、分母 $(z-w)\to(w-z)$、双和的分母 $(z-z_j)(w-z_k)\to(z-z_k)(w-z_j)$。这一步在推导 KZ 方程时要与 (6.12) 相减，正是“对称化”的来源，核对时可以逐项比对四个分母。]
由 (6.8) 与定理 3.26，能量-动量张量的关联函数有如下表达式。
\begin{align}
&\langle\Psi|T(z)|v\rangle(dz)^2\nonumber\\
&\quad=\frac{1}{g^*+\ell}\sum_{j,k}\frac{1}{(z-z_j)(z-z_k)}\cdot\sum_{a=1}^{\dim\mathfrak{g}}\langle\Psi|\rho_j(J^a)\rho_k(J^a)|v\rangle(dz)^2.\tag{6.14}
\end{align}
对 $\langle\Psi|\in\mathcal{V}^{\dagger}_{\vec\lambda}(\mathfrak{X})$ 与 $|v\rangle\in V_{\vec\lambda}$。置
\begin{equation}
\Omega_{jk}=\sum_{a=1}^{\dim\mathfrak{g}}\rho_j(J^a)\rho_k(J^a).\tag{6.15}
\end{equation}
以及
\[
\langle\psi|=j(\langle\Psi|)\in\Hom_{\mathfrak{g}}(V_{\vec\lambda},\C).
\]
于是有下面将要反复用到的表达式。
\begin{equation}
\langle\Psi|T(z)|v\rangle(dz)^2=\frac{1}{g^*+\ell}\sum_{j=1}^{\dim\mathfrak{g}}\sum_{k=1}^{\dim\mathfrak{g}}\frac{\langle\psi|\Omega_{jk}|v\rangle}{(z-z_j)(z-z_k)}(dz)^2.\tag{6.16}
\end{equation}
% [Tutor Note: (6.16) 的求和上限原文印作 $\dim\mathfrak{g}$，按上下文应为 $N$（$j,k$ 标号的是插入点，不是基指标 $a$；$a$ 已在 (6.15) 中缩并掉）。又 (6.14) 与 (6.16) 是同一个量的两种写法：前者用 $\rho_j(J^a)\rho_k(J^a)$，后者用 $\Omega_{jk}$ 打包，核对时用 (6.15) 直接代入即可。]
% [原文 p.136]
\subsection{KZ 方程}
本节我们研究 $g=0$ 时真空空间所满足的微分方程。固定正整数 $N\ge 3$，并置
\[
B_N=\C^N\setminus\bigcup_{i<j}\Delta_{ij},
\]
其中
\[
\Delta_{ij}=\{(z_1,\dots,z_N)\in\C^N\mid z_i=z_j\}.
\]
我们把 $\PP^1$ 视为 $\C\cup\{\infty\}$，并在下文固定 $\C$ 的坐标 $u$。设 $\mathfrak{F}=(\pi:\mathcal{C}\to\mathcal{B};s_1,\dots,s_N;\eta_1,\dots,\eta_N)$ 是一族 $N$-带点黎曼球面，其中
\begin{align*}
\mathcal{B}&=B_N,\qquad \mathcal{C}=B_N\times\PP^1,\\
\pi&:\mathcal{C}\to\mathcal{B}\quad\text{到第二个因子的投影},\\
s_j&:\mathcal{B}\to\mathcal{B}\times\PP^1,\\
(z_1,\dots,z_N)&\mapsto(z_1,\dots,z_N,(z_j:1)),
\end{align*}
而
\[
\eta_j:\widehat{\mathcal{O}}_{\mathcal{C}/s_j(\mathcal{B})}\to\mathcal{O}_{\mathcal{B}}[[\xi]]
\]
由
\[
\eta_j^{-1}(\xi)=\xi_j=u-z_j
\]
给出。由于坐标 $\eta_j$ 由 $s_j$ 决定（注意我们固定了 $\C$ 的坐标 $u$），族 $\mathfrak{F}$ 可视为 $g=0$ 的 $N$-带点稳定曲线族。置
\[
\Lambda=\{1,2,\dots,N,\infty\},
\]
\[
\zeta_\lambda=u-z_\lambda,\quad \lambda=1,2,\dots,N,\qquad \zeta_\infty=1/u,
\]
并引入开覆盖 $\{\mathcal{U}_\lambda\}_{\lambda\in\Lambda}$，其定义为
\[
\mathcal{U}_\lambda=\{(z_1,\dots,z_N;\zeta_\lambda)\mid (z_1,\dots,z_N)\in B_N,\ |\zeta_\lambda|<\epsilon\}
\]
（$\epsilon$ 为适当的正数）。
% [Tutor Note: $B_N$ 是 $N$ 个点的位形空间（去掉对角判别式），$N\ge 3$ 才使 $B_N$ 是 $\dim=N$ 的光滑仿射簇、且对应的带点曲线稳定；$\zeta_\infty=1/u$ 是 $\infty$ 处的局部坐标，$\mathcal{U}_\infty$ 就是“$\infty$ 附近”的坐标卡。核对：$u=\zeta_\lambda+z_\lambda$，故 $\partial u/\partial\zeta_\lambda=1$，下文用 $\mathrm{d}u$ 与 $\mathrm{d}\zeta_\lambda$ 互换时只出符号。]
% [原文 p.137]
\begin{lemma}[6.3]
族 $\mathfrak{F}$ 的 Kodaira--Spencer 映射
\[
\rho:\Theta_{\mathcal{B}}\to R^1\pi_*\mathcal{O}_{\mathcal{C}/\mathcal{B}}\Big(-\sum_{j=1}^{N}s_j(\mathcal{B})\Big)
\]
由
\[
\theta=\sum a_j(z)\frac{\partial}{\partial z_j}\mapsto\{[\rho_{\lambda\mu}]\}
\]
给出，其中在 $\mathcal{U}_\lambda\cap\mathcal{U}_\mu\ne\emptyset$、$\lambda\ne\mu$ 上，$\rho_{\lambda\mu}$ 由
\[
\rho_{\lambda\mu}=\begin{cases}
(a_\mu(z)-a_\lambda(z))\dfrac{\partial}{\partial\zeta_\lambda} & \text{若 } 1\le\lambda,\mu\le N,\\[2mm]
-a_\lambda(z)\dfrac{\partial}{\partial\zeta_\lambda} & \text{若 } \mu=\infty
\end{cases}
\]
给出。此外，$\operatorname{Ker}\rho$ 是秩 3 的自由 $\mathcal{O}_{\mathcal{B}}$-模，由整体截面
\[
\sum_{j=1}^{N}z_j^{k}\frac{\partial}{\partial z_j},\qquad k=0,1,2
\]
张成。
\end{lemma}
\noindent\textit{证明.}\quad 对 $\mathcal{U}_\lambda\cap\mathcal{U}_\mu\ne\emptyset$、$\lambda\ne\mu$ 且 $1\le\lambda,\mu\le N$，有
\[
\zeta_\lambda=\zeta_\mu+z_\mu-z_\lambda.
\]
于是由 §1.3 得
\[
\rho_{\lambda\mu}=(a_\mu(z)-a_\lambda(z))\frac{\partial}{\partial\zeta_\lambda}.
\]
另一方面，若 $\mu=\infty$ 且 $\mathcal{U}_\lambda\cap\mathcal{U}_\mu\ne\emptyset$，则
\[
\zeta_\lambda=1/\zeta_\infty-z_\lambda.
\]
于是有
\[
\rho_{\lambda\infty}=-a_\lambda(z)\frac{\partial}{\partial\zeta_\lambda}.
\]
这就证明了引理的第一部分。
再设 $\{[\rho_{\lambda\mu}]\}$ 为零。则存在
\[
\theta_\lambda\in\Gamma\Big(\mathcal{U}_\lambda,\pi_*\mathcal{O}_{\mathcal{C}/\mathcal{B}}\Big(-\sum s_j(\mathcal{B})\Big)\Big)
\]
使得
\begin{align}
\theta_\lambda&=\vartheta_\lambda(z,\zeta_\lambda)\frac{\partial}{\partial\zeta_\lambda},\nonumber\\
\vartheta_\lambda(z,0)&=0,\nonumber\\
\rho_{\lambda\mu}&=\theta_\mu-\theta_\lambda.\tag{6.17}
\end{align}
% [Tutor Note: (6.17) 的三行要一起读：$\theta_\lambda$ 是 $\mathcal{U}_\lambda$ 上的截面，$\vartheta_\lambda(z,0)=0$ 是它在插入点处的归一化，第三行说 $\rho$ 的 Čech 1-上链是 $\{\theta_\lambda\}$ 的上边缘——即 Kodaira--Spencer 类的消失条件。核对办法：把第三行代回 $\lambda\to\mu\to\nu$ 循环即见上链条件自动满足。]
因此，若 $1\le\lambda,\mu\le N$，则有
\[
a_\mu(z)\frac{\partial}{\partial\zeta_\mu}-\theta_\lambda=a_\lambda(z)-\frac{\partial}{\partial\zeta_\lambda}
\]
而若 $\mu=\infty$，则有
\[
\theta_\infty=a_\lambda(z)\frac{\partial}{\partial\zeta_\lambda}-\theta_\lambda.
\]
% [原文 p.138]
% [Tutor Note: 上面第一式的原文印刷有误：正确形式应为 $a_\mu(z)\partial/\partial\zeta_\mu-\theta_\mu=a_\lambda(z)\partial/\partial\zeta_\lambda-\theta_\lambda$（两边都是由同一整体向量场 $\{a_\lambda\partial/\partial\zeta_\lambda-\theta_\lambda\}$ 在各自坐标卡上的限制）。照排以便对照，见报告 ERRATA。]
于是
\[
\Big\{a_\lambda(z)\frac{\partial}{\partial\zeta_\lambda}-\theta_\lambda\Big\}
\]
是 $\mathcal{C}$ 上的整体向量场，因此具有形式
\[
(A(z)+B(z)u+C(z)u^2)\frac{\partial}{\partial u}.
\]
于是有
\begin{align*}
\vartheta_\lambda(z,\zeta_\lambda)&=a_\lambda(z)-\{A(z)+B(z)(\zeta_\lambda+z_\lambda)+C(z)(\zeta_\lambda+z_\lambda)^2\}\\
&=a_\lambda(z)-\{A(z)+B(z)z_\lambda+C(z)z_\lambda^2\}+O(\zeta_\lambda).
\end{align*}
因此，由 (6.17) 令 $\zeta_\lambda=0$，即得
\[
a_\lambda(z)=\{A(z)+B(z)z_\lambda+C(z)z_\lambda^2\}.
\]
\hfill$\square$
由正合列
\[
0\to\Theta_{\mathcal{C}/\mathcal{B}}\Big(-\sum s_j(\mathcal{B})\Big)\to\Theta_{\mathcal{C}/\mathcal{B}}\to\bigoplus_{j=1}^{N}\mathcal{O}_{\mathcal{B}}\frac{d}{d\xi_j}\to0
\]
得到正合列
\[
0\to\pi_*\Theta_{\mathcal{C}/\mathcal{B}}\to\bigoplus_{j=1}^{N}\mathcal{O}_{\mathcal{B}}\frac{d}{d\xi_j}\xrightarrow{\vartheta}R^1\pi_*\Theta_{\mathcal{C}/\mathcal{B}}\Big(-\sum s_j(\mathcal{B})\Big)\to0.
\]
注意 $\pi_*\Theta_{\mathcal{C}/\mathcal{B}}$ 是秩 3 的自由 $\mathcal{O}_{\mathcal{B}}$-模，我们可以取
\[
u^{k}\frac{\partial}{\partial u},\qquad k=0,1,2
\]
作为 $\mathcal{O}_{\mathcal{B}}$-模的基。对元素 $\vec\ell\in\bigoplus_{j=1}^{N}\mathcal{O}_{\mathcal{B}}\frac{d}{d\xi_j}$，像 $\vartheta(\vec\ell)$ 由 1-上链 $\{\vartheta_{\lambda\mu}\}$ 给出，其中
\[
\vartheta_{\lambda\mu}=\begin{cases}\ell_\mu-\ell_\lambda & \text{若 } 1\le\lambda,\mu\le N,\\ -\ell_\lambda & \text{若 } \mu=\infty.\end{cases}
\]
由引理 6.3 得到如下引理。
\begin{lemma}[6.4]
存在自然的 $\mathcal{O}_{\mathcal{B}}$-模同构
\[
\theta:\ \bigoplus_{j=1}^{N}\mathcal{O}_{\mathcal{B}}\frac{d}{d\xi_j}\longrightarrow\Theta_{\mathcal{B}},\qquad
\sum_{j=1}^{N}a_j(z)\frac{d}{d\xi_j}\longmapsto-\sum_{j=1}^{N}a_j(z)\frac{d}{dz_j},
\]
使得下图交换。
\[
\begin{array}{ccccccccc}
0 & \longrightarrow & \pi_*\Theta_{\mathcal{C}/\mathcal{B}} & \longrightarrow & \bigoplus_{j=1}^{N}\mathcal{O}_{\mathcal{B}}\dfrac{d}{d\xi_j} & \xrightarrow{\ \vartheta\ } & R^1\pi_*\Theta_{\mathcal{C}/\mathcal{B}}\big(-\sum s_j(\mathcal{B})\big) & \longrightarrow & 0\\[3mm]
 & & \big\downarrow & & \begin{array}{c}\big\downarrow\\[1mm] \theta\end{array} & & \Vert & & \\[3mm]
0 & \longrightarrow & \operatorname{Ker}\rho & \longrightarrow & \Theta_{\mathcal{B}} & \xrightarrow{\ \rho\ } & R^1\pi_*\Theta_{\mathcal{C}/\mathcal{B}}\big(-\sum s_j(\mathcal{B})\big) & \longrightarrow & 0
\end{array}
\]
% [Tutor Note: 引理 6.4 的负号是全域性的：由 $\partial/\partial z_j=-\partial/\partial\zeta_j$（在 $\lambda$ 卡上 $u=\zeta_\lambda+z_\lambda$ 而 $\zeta_j=u-z_j$）可知 $\sum a_j d/d\xi_j$ 与 $-\sum a_j\partial/\partial z_j$ 在同一个 $\Theta_{\mathcal{B}}$ 类里。核对办法：把上式作用在 $u$ 上得 $-a_j$，而 $d/d\xi_j$ 作用在 $u$ 上得 $+1$，符号正好抵消。]
%  另：该同构式右端原文印 $\frac{d}{dz_j}$（同式左端即 $\frac{d}{d\xi_j}$），
%  非 $\frac{\partial}{\partial z_j}$；译文照排原文（见 ERRATA）。
% [原文 p.139]
其中在 $\pi_*\Theta_{\mathcal{C}/\mathcal{B}}$ 上同构 $\theta$ 具有形式
\[
\theta\Big(u^{k}\frac{\partial}{\partial u}\Big)=-\sum_{j=1}^{N}z_j^{k}\frac{\partial}{\partial z_j}.
\]
\end{lemma}
\begin{definition}[6.5]
对元素 $\vec\ell=(\ell_1,\dots,\ell_N)\in\bigoplus_{j=1}^{N}\mathcal{O}_{\mathcal{B}}\frac{d}{d\xi_j}$，我们如下定义 $\mathcal{H}_{\vec\lambda}(\mathfrak{F})$ 上的作用 $D(\vec\ell)$：
\begin{equation}
D(\vec\ell)(F\otimes|\Phi\rangle)=\theta(\vec\ell)(F)\otimes|\Phi\rangle+F\otimes\Big(\sum_{j=1}^{N}\rho_j(T[\ell_j])|\Phi\rangle\Big).\tag{6.18}
\end{equation}
\end{definition}
% [Tutor Note: $D(\vec\ell)$ 就是 KZ 联络的协变导数：第一项 $\theta(\vec\ell)(F)$ 是底空间方向上的求导（曲线族在 $\mathcal{B}$ 上的形变），第二项 $\sum_j\rho_j(T[\ell_j])$ 是场在插入点处的作用。两项的“扭曲”体现在 $\theta$ 把 $d/d\xi_j$ 换成 $-\partial/\partial z_j$ 上，故推论 6.7 称之为“扭曲一阶微分算子”。]
\begin{proposition}[6.6]
对每个元素 $\vec\ell\in\bigoplus_{j=1}^{N}\mathcal{O}_{\mathcal{B}}\frac{d}{d\xi_j}$，有
\[
D(\vec\ell)(\hat{\mathfrak{g}}(\mathfrak{F})\mathcal{H}_{\vec\lambda}(\mathfrak{F}))\subset\hat{\mathfrak{g}}(\mathfrak{F})\mathcal{H}_{\vec\lambda}(\mathfrak{F}).
\]
\end{proposition}
\noindent\textit{证明.}\quad 层 $\pi_*\mathcal{O}_{\mathcal{B}}(*\sum s_j(\mathcal{B}))$ 由
\[
\frac{1}{(u-z_i)^m},\qquad i=1,2,\dots,N,\quad m=0,1,2,\dots
\]
张成。由命题 5.4 的证明（特别是 (4.18)），只需证明：对 $\vec\ell=(\ell_1,\dots,\ell_N)\in\bigoplus_{j=1}^{N}\mathcal{O}_{\mathcal{B}}\frac{d}{d\xi_j}$ 与 $h\in\pi_*\mathcal{O}_{\mathcal{C}}(*\sum s_j(\mathcal{B}))$，有
\[
(\theta(\vec\ell)-\ell_j)(h)\in\pi_*\mathcal{O}_{\mathcal{C}}\Big(*\sum s_j(\mathcal{B})\Big).
\]
为此只需证明：对 $A(z)\in\mathcal{O}_{\mathcal{B}}$ 有
\[
(\theta(\vec\ell)-\ell_j)\Big(\frac{A(z)}{(u-z_j)^m}\Big)\in\pi_*\mathcal{O}_{\mathcal{C}}\Big(*\sum s_j(\mathcal{B})\Big).
\]
置
\[
\ell_j=c_j(z)\frac{d}{d\xi_j},\qquad j=1,2,\dots,N.
\]
则由引理 6.4 有
\begin{align*}
\theta(\vec\ell)\Big(t_j\Big(\frac{A(z)}{(u-z_j)^m}\Big)\Big)&=\frac{-\sum_{j=1}^{N}c_j(z)\frac{\partial A(z)}{\partial z_j}}{(u-z_j)^m}+\frac{(mc_j(z)-c_i(z))A(z)}{(u-z_i)^{m+1}},\\
\ell_j\Big(\frac{A(z)}{(u-z_i)^m}\Big)&=\frac{-mc_j(z)A(z)}{(u-z_i)^{m+1}}.
\end{align*}
于是有
\[
(\theta(\vec\ell)-\ell_j)\Big(t_j\Big(\frac{A(z)}{(u-z_j)^m}\Big)\Big)=\frac{-\sum_{j=1}^{N}c_j(z)\frac{\partial A(z)}{\partial z_j}}{(u-z_j)^m}-\frac{mc_i(z)A(z)}{(u-z_i)^{m+1}}.
\]
右端与 $j$ 无关，并且定义了 $\pi_*\mathcal{O}_{\mathcal{C}}(*\sum s_j(\mathcal{B}))$ 的一个元素。\hfill$\square$
% [Tutor Note: 这里的 $t_j$ 是“往第 $j$ 个张量因子上插入”的记号（原文未定义，见命题 5.4 的证明）：$t_j(h)$ 把 $h\in\mathcal{O}_{\mathcal{C}}$ 放到第 $j$ 个位置。上面两式的指标 $i$ 与 $j$ 在原文中混用（第一式左端指标为 $j$、第二项分母却用 $i$），照排；只需注意结论是“右端与 $j$ 无关”，故这类混用不影响命题。核对办法：把 $m=0$ 代入，第一式退化为 $\theta(\vec\ell)(A(z))$ 的通常作用。]
\begin{corollary}[6.7]
对每个元素 $\vec\ell\in\bigoplus_{j=1}^{N}\mathcal{O}_{\mathcal{B}}\frac{d}{d\xi_j}$，$D(\vec\ell)$ 作为扭曲一阶微分算子作用在 $\mathcal{V}_{\vec\lambda}(\mathfrak{F})$ 与 $\mathcal{V}^{\dagger}_{\vec\lambda}(\mathfrak{F})$ 上。
\end{corollary}
% [原文 p.140]
\begin{lemma}[6.8]
向量空间
\[
H^0(\PP^1\times\PP^1,\omega_{\PP^1\times\PP^1}(2\Delta))
\]
是一维的，并由元素
\begin{equation}
\omega=\frac{dw\wedge dz}{(w-z)^2}\tag{6.19}
\end{equation}
张成。特别地，
\[
S_\omega\equiv0.
\]
\end{lemma}
由该引理，在我们的情形可以在 $\mathcal{V}_{\vec\lambda}(\mathfrak{F})$ 与 $\mathcal{V}^{\dagger}_{\vec\lambda}(\mathfrak{F})$ 上定义可积联络。下面把该联络显式写出。对任意 $\langle\Psi|\in\mathcal{V}^{\dagger}_{\vec\lambda}(\mathfrak{F})$、$v\in V_{\vec\lambda}$ 与 $\vec\ell\in\bigoplus_{j=1}^{N}\mathcal{O}_{\mathcal{B}}\frac{d}{d\xi_j}$，由 (4.15) 与 (4.17) 有
\begin{align*}
\langle\Psi|\{D(\vec\ell)|v\rangle\}&=\sum_{j=1}^{N}\operatorname{Res}_{\xi_j=0}\big(\ell_j(\xi_j)\langle\Psi|T(\xi_j)|v\rangle d\xi_j\big),\\
\{\langle\Psi|D(\vec\ell)\}|v\rangle&=-\sum_{j=1}^{N}\operatorname{Res}_{\xi_j=0}\big(\ell_j(\xi_j)\langle\Psi|T(\xi_j)|v\rangle d\xi_j\big)+\theta(\vec\ell)(\langle\Psi|v\rangle).
\end{align*}
注意，由 §3.5 的论证，$\langle\Psi|D(\vec\ell)=0$ 当且仅当对一切 $v\in V_{\vec\lambda}$ 有 $\{\langle\Psi|D(\vec\ell)\}|v\rangle=0$。因此由上述等式可推出：$\langle\Psi|D(\vec\ell)=0$ 当且仅当如下等式成立：
\[
\theta(\vec\ell)\langle\Psi|v\rangle=\sum_{j=1}^{N}\operatorname{Res}_{\xi_j=0}\big(\ell_j(\xi_j)\langle\Psi|T(\xi_j)|v\rangle d\xi_j\big).
\]
由引理 6.8、(5.5) 与 (6.14)，定义平坦截面的微分方程由下式给出：
\begin{align*}
\frac{\partial}{\partial z_j}\langle\Psi|v\rangle&=\sum_{j+1}^{N}\operatorname{Res}_{\xi_j=0}\big(\langle\Psi|T(\xi_j)|v\rangle d\xi_j\big)\\
&=\frac{1}{(g^*+\ell)}\sum_{j=1}^{N}\operatorname{Res}_{\xi_j=0}\Big(\sum_{i=1}^{N}\sum_{k=1}^{N}\frac{\langle\Psi|\Omega_{ik}|v\rangle}{(\xi_j+z_j-z_i)(\xi_j+z_j-z_k)}d\xi_j\Big)\\
&=\frac{1}{(g^*+\ell)}\sum_{k=1,\,k\ne j}^{N}\frac{\langle\Psi|\Omega_{jk}|v\rangle}{z_j-z_k}.
\end{align*}
这些方程就是 Knizhnik--Zamolodchikov 方程（[KZ]）。为简便起见，下文称它们为 KZ 方程。
% [Tutor Note: KZ 方程第一行的求和下限原文印作 $j+1$（显然应为 $j=1$），照排。第三行只留下 $k\ne j$ 的项：因为 $\mathrm{Res}_{\xi_j=0}$ 只挑出 $1/\xi_j$ 的系数，而在第二行中 $\xi_j+z_j-z_i$ 当 $i=j$ 时为 $\xi_j$、当 $k=j$ 时为 $\xi_j$，故 $i=k=j$ 的项给出 $1/\xi_j^2$（无留数），只有 $i=k\ne j$ 与 $i\ne k$ 的交叉项按 $1/\xi_j$ 取出，最终化成 $\Omega_{jk}/(z_j-z_k)$。核对办法：取 $N=1$ 时右端为零，与“单点球面无模”一致。]
